% Part: incompleteness
% Chapter: incompleteness-provability
% Section: introduction

\documentclass[../../../include/open-logic-section]{subfiles}

\begin{document}

\olfileid{inc}{inp}{int}

\olsection{مقدمه}

هیلبرت بر آن بود که دستگاهی از اصول موضوعه برای یک ساختار ریاضی، مانند
اعداد طبیعی، نارساست مگر آنکه بتوان همهٔ گزاره‌های صادق دربارهٔ آن
ساختار را در آن به دست آورد. این دیدگاه، همراه با علاقهٔ بعدی او به
دستگاه‌های صوری استنتاج، نشان می‌دهد که گمان می‌کرد باید تضمین کنیم
دستگاه‌های صوری‌ای که مثلاً برای استدلال دربارهٔ اعداد طبیعی به کار
می‌بریم نه‌تنها سازگار، بلکه \emph{کامل} نیز باشند؛ یعنی هر گزاره در
زبان دستگاه یا !!{derivable} باشد یا نقیض آن چنین باشد. قضیهٔ نخست
ناتمامیت گودل نشان می‌دهد که هیچ دستگاه اصول موضوعه‌ای از این دست وجود
ندارد: هیچ دستگاه صوری کامل و سازگاری برای حساب وجود ندارد که
!!{axiomatizable} باشد. در واقع، هیچ نظریهٔ ریاضیِ «به اندازهٔ کافی
نیرومند» و سازگاری که !!{axiomatizable} باشد کامل نیست.

هدف مهم‌تر هیلبرت، که کانون برنامهٔ او برای توجیه ریاضیات جدید
(«کلاسیک») بود، یافتن اثبات‌های متناهی‌گرایانهٔ سازگاری برای دستگاه‌های
صوریِ بازنمایندهٔ استدلال کلاسیک بود. ازاین‌رو، در نسبت با برنامهٔ
هیلبرت، قضیهٔ دوم ناتمامیت گودل ضربه‌ای بسیار سنگین‌تر به شمار می‌آمد.
قضیهٔ دوم ناتمامیت را، مانند قضیهٔ نخست، می‌توان به‌صورتی کمابیش مبهم
بیان کرد: هیچ نظریهٔ حسابِ سازگار و به اندازهٔ کافی نیرومندی نمی‌تواند سازگاری
خود را ثابت کند. در اینجا «به اندازهٔ کافی نیرومند» باید اندکی بیش از
$\Th{Q}$ را در بر گیرد.

اندیشهٔ اصلی اثبات نخست گودل را می‌توان در پارادوکس اپیمنیدس یافت.
اپیمنیدس، که اهل کرت بود، ادعا کرد همهٔ کرتیان دروغ‌گویند؛ صورت مستقیم‌تر
پارادوکس گزارهٔ «این جمله کاذب است» است. گودل، در اصل با جایگزین‌کردن
صدق با !!{derivability}، توانست !!a{sentence}ی را صوری کند که به شیوه‌ای
غیرمستقیم می‌گوید خودش !!{derivable} نیست. اگر آن !!{sentence}
!!{derivable} می‌بود، نظریه ناسازگار می‌شد. گودل همچنین نشان داد نقیض آن
!!{sentence} نیز از دستگاه اصول موضوعهٔ مورد نظر او !!{derivable}
نیست. (برای این بخش دوم، گودل ناچار بود فرض کند نظریهٔ~$\Th{T}$
«$\omega$-سازگار» است. $\omega$-سازگاری با سازگاری مرتبط، اما خاصیتی
قوی‌تر است.\footnote{یعنی هر نظریهٔ $\omega$-سازگار سازگار است، اما عکس
آن برقرار نیست.} چند سال پس از گودل، راسر نشان داد فرض سازگاری سادهٔ
$\Th{T}$ کافی است.)

نخستین چالش، فهم این نکته است که چگونه می‌توان !!a{sentence}ی ساخت که
به خودش ارجاع دهد. برای هر !!{formula}~$!A$ در زبان~$\Th{Q}$، بگذارید
$\gn{!A}$ عددنمای متناظر با~$\Gn{!A}$ باشد. معنای این قرارداد را در
نظر بگیرید: $!A$ !!a{formula} در زبان~$\Th{Q}$ است، $\Gn{!A}$ عددی
طبیعی است و~$\gn{!A}$ \emph{ترمی} در زبان~$\Th{Q}$ است. پس هر
!!{formula}~$!A$ در زبان~$\Th{Q}$ \emph{نامی} چون~$\gn{!A}$ دارد که
ترمی در زبان~$\Th{Q}$ است؛ این امر چارچوبی مفهومی فراهم می‌کند که در آن
!!{formula}s در زبان~$\Th{Q}$ می‌توانند دربارهٔ !!{formula}sی دیگر
«سخن بگویند». لم زیر به لم نقطهٔ ثابت معروف است.

\begin{lem}
فرض کنید~$\Th{T}$ نظریه‌ای دلخواه و گسترشی از~$\Th{Q}$ باشد، و
$!B(x)$ هر !!{formula}ی باشد که تنها متغیر آزادش~$x$ است. آنگاه
!!a{sentence}ی چون~$!A$ وجود دارد که
$\Th{T} \Proves!A \liff !B(\gn{!A})$.
\end{lem}

لم می‌گوید برای هر خاصیت~$!B(x)$، !!a{sentence}ی چون~$!A$ وجود دارد که
می‌گوید «$!B(x)$ دربارهٔ من صادق است»، و~$\Th{T}$ این را «می‌داند».

چگونه می‌توانیم چنین !!a{sentence}ی بسازیم؟ صورت زیر از پارادوکس
اپیمنیدس را که از کواین است در نظر بگیرید:
\begin{quote}
«اگر نقل‌قول خودش پیش از آن بیاید کذب می‌دهد» اگر نقل‌قول خودش پیش از
آن بیاید کذب می‌دهد.
\end{quote}
این جمله مستقیماً خودارجاع نیست؛ فقط دربارهٔ اشیای نحوی میان گیومه‌ها
ادعایی می‌کند و از این حیث هم‌ردیف جمله‌هایی از این دست است:
\begin{enumerate}
\item «رابرت» نام زیبایی است.
\item «دویدم.» جملهٔ کوتاهی است.
\item «سه واژه دارد» سه واژه دارد.
\end{enumerate}
اما اگر عبارت «اگر نقل‌قول خودش پیش از آن بیاید کذب می‌دهد» را برداریم
و صورت نقل‌قول‌شدهٔ خودش را پیش از آن قرار دهیم چه رخ می‌دهد؟ همان جملهٔ
اصلی را به دست می‌آوریم!{} خلاصه اینکه جمله ادعا می‌کند خودش کاذب است.

\end{document}
